% Standalone problem document, generated from the adjacent README.md.
% Compile with XeLaTeX. Mathematical content is maintained in Markdown.
\documentclass[11pt,a4paper]{article}
\usepackage[fontsize=10.5pt]{scrextend}
\usepackage[margin=22mm,headheight=15pt,headsep=9mm,footskip=11mm]{geometry}
\usepackage{fontspec}
\setmainfont{texgyrepagella}[Extension=.otf,UprightFont=*-regular,BoldFont=*-bold,ItalicFont=*-italic,BoldItalicFont=*-bolditalic]
\setsansfont{texgyreheros}[Extension=.otf,UprightFont=*-regular,BoldFont=*-bold,ItalicFont=*-italic,BoldItalicFont=*-bolditalic]
\setmonofont{texgyrecursor}[Extension=.otf,UprightFont=*-regular,BoldFont=*-bold,ItalicFont=*-italic,BoldItalicFont=*-bolditalic,Scale=MatchLowercase]
\usepackage{amsmath,amssymb}
\usepackage{unicode-math}
\setmathfont{texgyrepagella-math.otf}
\usepackage{xcolor,microtype,enumitem,needspace,fancyhdr}
\usepackage{bookmark}
\definecolor{ink}{HTML}{17344B}
\definecolor{accent}{HTML}{197D80}
\definecolor{muted}{HTML}{596773}
\hypersetup{colorlinks=true,linkcolor=accent,urlcolor=accent,pdftitle={PF-02 - Connectedness
of minimal positive semidefinite factorization
orbits},pdfauthor={Open Problems in Numerical Linear Algebra}}
\urlstyle{same}
\setlength{\parindent}{0pt}
\setlength{\parskip}{5pt plus 1pt minus 1pt}
\linespread{1.04}
\setlength{\emergencystretch}{3em}
\widowpenalty=10000
\clubpenalty=10000
\displaywidowpenalty=10000
\allowdisplaybreaks[1]
\setlist{nosep,leftmargin=1.5em}
\providecommand{\tightlist}{\setlength{\itemsep}{0pt}\setlength{\parskip}{0pt}}
\providecommand{\pandocbounded}[1]{#1}
\pagestyle{fancy}
\fancyhf{}
\fancyhead[L]{\sffamily\footnotesize\color{muted}OPEN PROBLEMS IN NUMERICAL LINEAR ALGEBRA}
\fancyhead[R]{\sffamily\bfseries\footnotesize\color{accent}PF-02}
\fancyfoot[L]{\parbox[b]{0.9\textwidth}{\sffamily\fontsize{8}{10}\selectfont\color{muted}Editorial ratings; a bounded search cannot certify that a problem remains unsolved.\\Literature check: 2026-09-08}}
\fancyfoot[R]{\sffamily\footnotesize\color{muted}\thepage}
\renewcommand{\headrulewidth}{0.3pt}
\renewcommand{\headrule}{\hbox to\headwidth{\color{accent}\leaders\hrule height \headrulewidth\hfill}}
\makeatletter
\renewcommand{\section}{\Needspace{5\baselineskip}\@startsection{section}{1}{0pt}{12pt plus 2pt minus 2pt}{4pt}{\sffamily\bfseries\large\color{ink}}}
\renewcommand{\subsection}{\Needspace{4\baselineskip}\@startsection{subsection}{2}{0pt}{9pt plus 1pt minus 1pt}{3pt}{\sffamily\bfseries\normalsize\color{ink}}}
\makeatother
\setcounter{secnumdepth}{0}
\begin{document}
{\sffamily\small\color{muted}Nonnegative and positive
factorizations\par}
\vspace{3pt}
{\raggedright\hyphenpenalty=10000\exhyphenpenalty=10000\sffamily\bfseries\fontsize{21}{24}\selectfont\color{ink}Connectedness
of minimal positive semidefinite factorization orbits\par}
\vspace{7pt}
{\sffamily\small\textbf{Difficulty:} challenging\quad\textbf{Importance:} interesting
to the community\par}
\vspace{4pt}
{\color{accent}\hrule height 0.8pt}
\vspace{6pt}
\textbf{Status:} open; checked 2026-09-08\\
\textbf{Area:} geometry of constrained matrix factorizations

\subsection{Context and notation}\label{context-and-notation}

Write \(\mathbb S_+^k\) for the cone of real symmetric positive
semidefinite \(k\times k\) matrices. For an entrywise nonnegative matrix
\(M\in\mathbb R_+^{p\times q}\), its \textbf{real positive semidefinite
rank} is

\[
\operatorname{rank}_{\rm psd}(M)
=\min\{k\ge1:\ \exists A_1,\ldots,A_p,B_1,\ldots,B_q\in\mathbb S_+^k,\quad
M_{ij}=\operatorname{tr}(A_iB_j)\ \text{for all }i,j\}.
\]

This definition concerns a family of matrix factors indexed by the rows
and columns of \(M\).

\subsection{Problem statement}\label{problem-statement}

Let \(k\ge3\) and \(p,q\ge1\) be integers. Suppose
\(M\in\mathbb R_+^{p\times q}\) satisfies

\[
\operatorname{rank}(M)=\frac{k(k+1)}2,
\qquad \operatorname{rank}_{\rm psd}(M)=k.
\]

Let \(\mathcal F_k(M)\) be the set of all tuples
\((A_1,\ldots,A_p,B_1,\ldots,B_q)\) in \((\mathbb S_+^k)^{p+q}\)
satisfying \(\operatorname{tr}(A_iB_j)=M_{ij}\) for every \(i,j\). Give
it the Euclidean subspace topology. Identify tuples under the changes of
basis

\[
A_i\longmapsto S^\mathsf T A_iS,\qquad
B_j\longmapsto S^{-1}B_jS^{-\mathsf T},
\qquad S\in GL(k,\mathbb R).
\]

Give the resulting orbit space the quotient topology.

\subsubsection{Question}\label{question}

Is \(\mathcal F_k(M)/GL(k,\mathbb R)\) connected for every \(M\)
satisfying these hypotheses?

The ordinary-rank hypothesis is part of the problem. The case \(k=2\) is
known to be connected. The question asks whether optimal factors can
belong to separated families after changes of basis are identified.

\subsection{References}\label{references}

Fawzi, Gouveia, Parrilo, Robinson, and Thomas,
\href{https://arxiv.org/html/1407.4095}{\emph{Positive semidefinite
rank}}, §9.2, Problem 9.4. Richard Z. Robinson,
\href{https://digital.lib.washington.edu/server/api/core/bitstreams/4e9d6133-5d14-4071-9905-70bd7dfd530e/content}{\emph{The
Positive Semidefinite Rank of Matrices and Polytopes}}, University of
Washington dissertation (2015), Chapter 7, especially Proposition 7.0.8.

\subsection{Status evidence}\label{status-evidence}

Searches for ``psd factorizations connected'', ``positive semidefinite
factorization space connected'', and ``psd factorization orbits
connected 2026'' found the original question and the known \(k=2\)
result, but no resolution under the displayed rank conditions.
Universality or disconnectedness results for nonnegative vector
factorizations do not by themselves answer this question. No recent
explicit open-status confirmation was located.
\end{document}
